\subsection{Hitos de recepción y pagos de implementación}
Los doce hitos E-25 conservan sus meses y condiciones de devengo. La Etapa 1 suma 60\,\% y la Etapa 2 40\,\% del valor de implementación. Una fecha programada no constituye aprobación ni autoriza pago sin su entregable recibido \parencite[Formulario~E-25; p.~74]{bases-admin}.
La habilitación H3 comprende además DR conforme a RT-04.01 de las Bases Técnicas Transversales, aunque la descripción de E-25 enumere DEV, QA, PREPROD y PROD. La EDT incluye una ficha por cada uno de los cinco ambientes \parencite[RT-04.01; pp.~10--11]{bases-transversales}.
La Tabla~\ref{tab:sd7-e25} relaciona los doce hitos con mes, porcentaje y condición de recepción.
\begin{longtable}{L{12mm}R{12mm}R{14mm}L{110mm}}
\caption{Correspondencia contractual con E-25}\label{tab:sd7-e25}\\\hline
 \EncabezadoTabla{Hito} & \EncabezadoTabla{Mes} & \EncabezadoTabla{Pago} & \EncabezadoTabla{Condición de devengo}\\\hline\endfirsthead
\hline \EncabezadoTabla{Hito} & \EncabezadoTabla{Mes} & \EncabezadoTabla{Pago} & \EncabezadoTabla{Condición de devengo}\\\hline\endhead
H1 & 2 & 8\,\% & Alcance y matriz de trazabilidad aprobados \\\hline
H2 & 4 & 7\,\% & Arquitectura, seguridad y modelo de datos aprobados\\\hline
H3 & 6 & 10\,\% & Infraestructura híbrida y cuatro ambientes con observabilidad \\\hline
H4 & 10 & 10\,\% & Software E1 con QA superado y evidencia de cobertura\\\hline
H5 & 12 & 10\,\% & E1 certificada: UAT, carga, resiliencia y seguridad ofensiva \\\hline
H6 & 13 & 5\,\% & Inicio de marcha blanca E1; retorno activo y usuarios capacitados\\\hline
H7 & 16 & 10\,\% & Producción E1 y cierre conforme de marcha blanca \\\hline
H8 & 14 & 5\,\% & Alcance E2 y diseño detallado aprobados\\\hline
H9 & 17 & 10\,\% & Software E2 con QA superado \\\hline
H10 & 18 & 10\,\% & Certificación E2 y cierre de desarrollo\\\hline
H11 & 19 & 5\,\% & Inicio de marcha blanca E2 con E1 en producción \\\hline
H12 & 21 & 10\,\% & Producción E2 y aceptación final de implementación\\\hline
\end{longtable}\FuenteTabla{Bases Administrativas, Formulario E-25. Porcentajes sobre el valor total de implementación.}


La operación comprende 36 períodos, meses 21--56, con pago al mes siguiente vencido sujeto a los descuentos de SLA aplicables. El último período de servicio es el 56 y su pago corresponde al mes siguiente; ese pago no agrega un mes de operación. No se anticiparán ni prorratearán estos importes dentro de implementación. Los porcentajes E-25 se convertirán en montos al consolidar la valoración económica, conservando la separación entre ambas fases.

\subsection{Gantt contractual de 56 meses}
La carta muestra los períodos de trabajo de las doce cuentas, marchas blancas, cobertura inicial y Operación, junto con los doce hitos E-25. IN2 conserva junio de 2029 como primera invernada; las demás innovaciones siguen SD13. Las barras agregadas no autorizan intervenciones durante congelamientos, eventos o marejadas ni sustituyen las precedencias de la red.~\parencite[art.~17, pp.~12--13; Formulario E-25, p.~74]{bases-admin}
\clearpage\begin{landscape}
La Figura~\ref{fig:sd7-gantt-1} presenta la programación agregada de los meses 1--28.
La carta agrega los intervalos calculados por cuenta; las fases contractuales se muestran en filas separadas. Las barras no autorizan cambios durante un congelamiento ni acreditan recepciones ejecutadas.
\begin{figure}[H]\centering\begin{tikzpicture}[x=0.52cm,y=0.55cm,font=\fontsize{9}{11}\selectfont]
\node at (0.5,.6) {1};\draw[SynLine] (0,.2)--(0,-16.5);
\node at (1.5,.6) {2};\draw[SynLine] (1,.2)--(1,-16.5);
\node at (2.5,.6) {3};\draw[SynLine] (2,.2)--(2,-16.5);
\node at (3.5,.6) {4};\draw[SynLine] (3,.2)--(3,-16.5);
\node at (4.5,.6) {5};\draw[SynLine] (4,.2)--(4,-16.5);
\node at (5.5,.6) {6};\draw[SynLine] (5,.2)--(5,-16.5);
\node at (6.5,.6) {7};\draw[SynLine] (6,.2)--(6,-16.5);
\node at (7.5,.6) {8};\draw[SynLine] (7,.2)--(7,-16.5);
\node at (8.5,.6) {9};\draw[SynLine] (8,.2)--(8,-16.5);
\node at (9.5,.6) {10};\draw[SynLine] (9,.2)--(9,-16.5);
\node at (10.5,.6) {11};\draw[SynLine] (10,.2)--(10,-16.5);
\node at (11.5,.6) {12};\draw[SynLine] (11,.2)--(11,-16.5);
\node at (12.5,.6) {13};\draw[SynLine] (12,.2)--(12,-16.5);
\node at (13.5,.6) {14};\draw[SynLine] (13,.2)--(13,-16.5);
\node at (14.5,.6) {15};\draw[SynLine] (14,.2)--(14,-16.5);
\node at (15.5,.6) {16};\draw[SynLine] (15,.2)--(15,-16.5);
\node at (16.5,.6) {17};\draw[SynLine] (16,.2)--(16,-16.5);
\node at (17.5,.6) {18};\draw[SynLine] (17,.2)--(17,-16.5);
\node at (18.5,.6) {19};\draw[SynLine] (18,.2)--(18,-16.5);
\node at (19.5,.6) {20};\draw[SynLine] (19,.2)--(19,-16.5);
\node at (20.5,.6) {21};\draw[SynLine] (20,.2)--(20,-16.5);
\node at (21.5,.6) {22};\draw[SynLine] (21,.2)--(21,-16.5);
\node at (22.5,.6) {23};\draw[SynLine] (22,.2)--(22,-16.5);
\node at (23.5,.6) {24};\draw[SynLine] (23,.2)--(23,-16.5);
\node at (24.5,.6) {25};\draw[SynLine] (24,.2)--(24,-16.5);
\node at (25.5,.6) {26};\draw[SynLine] (25,.2)--(25,-16.5);
\node at (26.5,.6) {27};\draw[SynLine] (26,.2)--(26,-16.5);
\node at (27.5,.6) {28};\draw[SynLine] (27,.2)--(27,-16.5);
\node[anchor=east,align=right,text width=4.3cm] at (-.2,-1) {1.1 Gobierno y contrato};
\fill[SynBlue!55] (0,-1.22) rectangle (1,-0.78);
\fill[SynBlue!55] (1,-1.22) rectangle (2,-0.78);
\fill[SynBlue!55] (2,-1.22) rectangle (3,-0.78);
\fill[SynBlue!55] (3,-1.22) rectangle (4,-0.78);
\fill[SynBlue!55] (4,-1.22) rectangle (5,-0.78);
\fill[SynBlue!55] (5,-1.22) rectangle (6,-0.78);
\fill[SynBlue!55] (6,-1.22) rectangle (7,-0.78);
\fill[SynBlue!55] (7,-1.22) rectangle (8,-0.78);
\fill[SynBlue!55] (8,-1.22) rectangle (9,-0.78);
\fill[SynBlue!55] (9,-1.22) rectangle (10,-0.78);
\fill[SynBlue!55] (10,-1.22) rectangle (11,-0.78);
\fill[SynBlue!55] (11,-1.22) rectangle (12,-0.78);
\fill[SynBlue!55] (12,-1.22) rectangle (13,-0.78);
\fill[SynBlue!55] (13,-1.22) rectangle (14,-0.78);
\fill[SynBlue!55] (14,-1.22) rectangle (15,-0.78);
\fill[SynBlue!55] (15,-1.22) rectangle (16,-0.78);
\fill[SynBlue!55] (16,-1.22) rectangle (17,-0.78);
\fill[SynBlue!55] (17,-1.22) rectangle (18,-0.78);
\fill[SynBlue!55] (18,-1.22) rectangle (19,-0.78);
\fill[SynBlue!55] (19,-1.22) rectangle (20,-0.78);
\node[anchor=east,align=right,text width=4.3cm] at (-.2,-2) {1.2 Plataforma y activos};
\fill[SynBlue!55] (2,-2.22) rectangle (3,-1.78);
\fill[SynBlue!55] (3,-2.22) rectangle (4,-1.78);
\fill[SynBlue!55] (4,-2.22) rectangle (5,-1.78);
\node[anchor=east,align=right,text width=4.3cm] at (-.2,-3) {1.3 Seguridad y auditoría};
\fill[SynBlue!55] (1,-3.22) rectangle (2,-2.78);
\fill[SynBlue!55] (2,-3.22) rectangle (3,-2.78);
\fill[SynBlue!55] (3,-3.22) rectangle (4,-2.78);
\fill[SynBlue!55] (4,-3.22) rectangle (5,-2.78);
\fill[SynBlue!55] (7,-3.22) rectangle (8,-2.78);
\fill[SynBlue!55] (8,-3.22) rectangle (9,-2.78);
\node[anchor=east,align=right,text width=4.3cm] at (-.2,-4) {1.4 Núcleo operacional};
\fill[SynBlue!55] (7,-4.22) rectangle (8,-3.78);
\fill[SynBlue!55] (8,-4.22) rectangle (9,-3.78);
\fill[SynBlue!55] (1,-4.22) rectangle (2,-3.78);
\fill[SynBlue!55] (2,-4.22) rectangle (3,-3.78);
\node[anchor=east,align=right,text width=4.3cm] at (-.2,-5) {1.5 Núcleo de servicios};
\fill[SynBlue!55] (1,-5.22) rectangle (2,-4.78);
\fill[SynBlue!55] (2,-5.22) rectangle (3,-4.78);
\fill[SynBlue!55] (7,-5.22) rectangle (8,-4.78);
\fill[SynBlue!55] (8,-5.22) rectangle (9,-4.78);
\fill[SynBlue!55] (12,-5.22) rectangle (13,-4.78);
\fill[SynBlue!55] (13,-5.22) rectangle (14,-4.78);
\fill[SynBlue!55] (14,-5.22) rectangle (15,-4.78);
\node[anchor=east,align=right,text width=4.3cm] at (-.2,-6) {1.6 Núcleo administrativo};
\fill[SynBlue!55] (1,-6.22) rectangle (2,-5.78);
\fill[SynBlue!55] (2,-6.22) rectangle (3,-5.78);
\fill[SynBlue!55] (7,-6.22) rectangle (8,-5.78);
\fill[SynBlue!55] (8,-6.22) rectangle (9,-5.78);
\fill[SynBlue!55] (12,-6.22) rectangle (13,-5.78);
\fill[SynBlue!55] (13,-6.22) rectangle (14,-5.78);
\fill[SynBlue!55] (14,-6.22) rectangle (15,-5.78);
\fill[SynBlue!55] (15,-6.22) rectangle (16,-5.78);
\node[anchor=east,align=right,text width=4.3cm] at (-.2,-7) {1.7 Canales e integración};
\fill[SynBlue!55] (3,-7.22) rectangle (4,-6.78);
\fill[SynBlue!55] (4,-7.22) rectangle (5,-6.78);
\fill[SynBlue!55] (5,-7.22) rectangle (6,-6.78);
\fill[SynBlue!55] (7,-7.22) rectangle (8,-6.78);
\fill[SynBlue!55] (8,-7.22) rectangle (9,-6.78);
\fill[SynBlue!55] (13,-7.22) rectangle (14,-6.78);
\node[anchor=east,align=right,text width=4.3cm] at (-.2,-8) {1.8 Migración e historia};
\fill[SynBlue!55] (0,-8.22) rectangle (1,-7.78);
\fill[SynBlue!55] (1,-8.22) rectangle (2,-7.78);
\fill[SynBlue!55] (2,-8.22) rectangle (3,-7.78);
\fill[SynBlue!55] (6,-8.22) rectangle (7,-7.78);
\fill[SynBlue!55] (8,-8.22) rectangle (9,-7.78);
\fill[SynBlue!55] (13,-8.22) rectangle (14,-7.78);
\node[anchor=east,align=right,text width=4.3cm] at (-.2,-9) {1.9 Innovaciones y pilotos};
\fill[SynBlue!55] (1,-9.22) rectangle (2,-8.78);
\fill[SynBlue!55] (2,-9.22) rectangle (3,-8.78);
\fill[SynBlue!55] (3,-9.22) rectangle (4,-8.78);
\fill[SynBlue!55] (12,-9.22) rectangle (13,-8.78);
\fill[SynBlue!55] (13,-9.22) rectangle (14,-8.78);
\fill[SynBlue!55] (14,-9.22) rectangle (15,-8.78);
\fill[SynBlue!55] (15,-9.22) rectangle (16,-8.78);
\fill[SynBlue!55] (16,-9.22) rectangle (17,-8.78);
\fill[SynBlue!55] (17,-9.22) rectangle (18,-8.78);
\fill[SynBlue!55] (18,-9.22) rectangle (19,-8.78);
\fill[SynBlue!55] (19,-9.22) rectangle (20,-8.78);
\fill[SynBlue!55] (20,-9.22) rectangle (21,-8.78);
\node[anchor=east,align=right,text width=4.3cm] at (-.2,-10) {1.10 Pruebas y correcciones};
\fill[SynBlue!55] (8,-10.22) rectangle (9,-9.78);
\fill[SynBlue!55] (9,-10.22) rectangle (10,-9.78);
\fill[SynBlue!55] (10,-10.22) rectangle (11,-9.78);
\fill[SynBlue!55] (11,-10.22) rectangle (12,-9.78);
\fill[SynBlue!55] (13,-10.22) rectangle (14,-9.78);
\fill[SynBlue!55] (14,-10.22) rectangle (15,-9.78);
\fill[SynBlue!55] (15,-10.22) rectangle (16,-9.78);
\fill[SynBlue!55] (16,-10.22) rectangle (17,-9.78);
\node[anchor=east,align=right,text width=4.3cm] at (-.2,-11) {1.11 Formación e implantación};
\fill[SynBlue!55] (10,-11.22) rectangle (11,-10.78);
\fill[SynBlue!55] (11,-11.22) rectangle (12,-10.78);
\fill[SynBlue!55] (12,-11.22) rectangle (13,-10.78);
\fill[SynBlue!55] (13,-11.22) rectangle (14,-10.78);
\fill[SynBlue!55] (14,-11.22) rectangle (15,-10.78);
\fill[SynBlue!55] (15,-11.22) rectangle (16,-10.78);
\fill[SynBlue!55] (16,-11.22) rectangle (17,-10.78);
\fill[SynBlue!55] (17,-11.22) rectangle (18,-10.78);
\fill[SynBlue!55] (18,-11.22) rectangle (19,-10.78);
\fill[SynBlue!55] (19,-11.22) rectangle (20,-10.78);
\fill[SynBlue!55] (20,-11.22) rectangle (21,-10.78);
\node[anchor=east,align=right,text width=4.3cm] at (-.2,-12) {1.12 Servicio y salida};
\fill[SynBlue!55] (17,-12.22) rectangle (18,-11.78);
\fill[SynBlue!55] (19,-12.22) rectangle (20,-11.78);
\fill[SynBlue!55] (20,-12.22) rectangle (21,-11.78);
\fill[SynBlue!55] (21,-12.22) rectangle (22,-11.78);
\fill[SynBlue!55] (22,-12.22) rectangle (23,-11.78);
\fill[SynBlue!55] (23,-12.22) rectangle (24,-11.78);
\fill[SynBlue!55] (24,-12.22) rectangle (25,-11.78);
\fill[SynBlue!55] (25,-12.22) rectangle (26,-11.78);
\fill[SynBlue!55] (26,-12.22) rectangle (27,-11.78);
\fill[SynBlue!55] (27,-12.22) rectangle (28,-11.78);
\node[anchor=east,text width=4.3cm,align=right] at (-.2,-13) {Marcha blanca E1};
\fill[SynTeal!65] (12,-13.22) rectangle (15,-12.78);
\node[anchor=east,text width=4.3cm,align=right] at (-.2,-14) {Marcha blanca E2};
\fill[SynTeal!65] (18,-14.22) rectangle (20,-13.78);
\node[anchor=east,text width=4.3cm,align=right] at (-.2,-15) {Operación ambos alcances};
\fill[SynNavyLight!65] (20,-15.22) rectangle (28,-14.78);
\node[anchor=east,text width=4.3cm,align=right] at (-.2,-16) {Cobertura inicial E1/E2};
\fill[SynNavyLight!40] (12,-16.22) rectangle (20,-15.78);
\node[rotate=90,anchor=west,text=SynNavy] at (1.5,-17.2) {H1};
\node[rotate=90,anchor=west,text=SynNavy] at (3.5,-17.2) {H2};
\node[rotate=90,anchor=west,text=SynNavy] at (5.5,-17.2) {H3};
\node[rotate=90,anchor=west,text=SynNavy] at (9.5,-17.2) {H4};
\node[rotate=90,anchor=west,text=SynNavy] at (11.5,-17.2) {H5};
\node[rotate=90,anchor=west,text=SynNavy] at (12.5,-17.2) {H6};
\node[rotate=90,anchor=west,text=SynNavy] at (15.5,-17.2) {H7};
\node[rotate=90,anchor=west,text=SynNavy] at (13.5,-17.2) {H8};
\node[rotate=90,anchor=west,text=SynNavy] at (16.5,-17.2) {H9};
\node[rotate=90,anchor=west,text=SynNavy] at (17.5,-17.2) {H10};
\node[rotate=90,anchor=west,text=SynNavy] at (18.5,-17.2) {H11};
\node[rotate=90,anchor=west,text=SynNavy] at (20.5,-17.2) {H12};
\end{tikzpicture}\caption{Carta Gantt agregada de los meses 1--28}\label{fig:sd7-gantt-1}
\FuenteFigura{Intervalos de red y asignaciones con esfuerzo UCP seleccionado; acta de marcha blanca E2 y corte requeridos al inicio de 21, bajo el supuesto adoptado de recepción continua y firma expresa programada.}
\end{figure}
Las filas permiten contrastar la continuidad de las cuentas con las ventanas contractuales. Los hitos identifican recepciones; la extensión de una barra no modifica las precedencias ni habilita trabajo durante una restricción operacional.
\end{landscape}\clearpage
\clearpage\begin{landscape}
La Figura~\ref{fig:sd7-gantt-2} presenta la programación agregada de los meses 29--56.
La carta agrega los intervalos calculados por cuenta; las fases contractuales se muestran en filas separadas. Las barras no autorizan cambios durante un congelamiento ni acreditan recepciones ejecutadas.
\begin{figure}[H]\centering\begin{tikzpicture}[x=0.52cm,y=0.55cm,font=\fontsize{9}{11}\selectfont]
\node at (0.5,.6) {29};\draw[SynLine] (0,.2)--(0,-16.5);
\node at (1.5,.6) {30};\draw[SynLine] (1,.2)--(1,-16.5);
\node at (2.5,.6) {31};\draw[SynLine] (2,.2)--(2,-16.5);
\node at (3.5,.6) {32};\draw[SynLine] (3,.2)--(3,-16.5);
\node at (4.5,.6) {33};\draw[SynLine] (4,.2)--(4,-16.5);
\node at (5.5,.6) {34};\draw[SynLine] (5,.2)--(5,-16.5);
\node at (6.5,.6) {35};\draw[SynLine] (6,.2)--(6,-16.5);
\node at (7.5,.6) {36};\draw[SynLine] (7,.2)--(7,-16.5);
\node at (8.5,.6) {37};\draw[SynLine] (8,.2)--(8,-16.5);
\node at (9.5,.6) {38};\draw[SynLine] (9,.2)--(9,-16.5);
\node at (10.5,.6) {39};\draw[SynLine] (10,.2)--(10,-16.5);
\node at (11.5,.6) {40};\draw[SynLine] (11,.2)--(11,-16.5);
\node at (12.5,.6) {41};\draw[SynLine] (12,.2)--(12,-16.5);
\node at (13.5,.6) {42};\draw[SynLine] (13,.2)--(13,-16.5);
\node at (14.5,.6) {43};\draw[SynLine] (14,.2)--(14,-16.5);
\node at (15.5,.6) {44};\draw[SynLine] (15,.2)--(15,-16.5);
\node at (16.5,.6) {45};\draw[SynLine] (16,.2)--(16,-16.5);
\node at (17.5,.6) {46};\draw[SynLine] (17,.2)--(17,-16.5);
\node at (18.5,.6) {47};\draw[SynLine] (18,.2)--(18,-16.5);
\node at (19.5,.6) {48};\draw[SynLine] (19,.2)--(19,-16.5);
\node at (20.5,.6) {49};\draw[SynLine] (20,.2)--(20,-16.5);
\node at (21.5,.6) {50};\draw[SynLine] (21,.2)--(21,-16.5);
\node at (22.5,.6) {51};\draw[SynLine] (22,.2)--(22,-16.5);
\node at (23.5,.6) {52};\draw[SynLine] (23,.2)--(23,-16.5);
\node at (24.5,.6) {53};\draw[SynLine] (24,.2)--(24,-16.5);
\node at (25.5,.6) {54};\draw[SynLine] (25,.2)--(25,-16.5);
\node at (26.5,.6) {55};\draw[SynLine] (26,.2)--(26,-16.5);
\node at (27.5,.6) {56};\draw[SynLine] (27,.2)--(27,-16.5);
\node[anchor=east,align=right,text width=4.3cm] at (-.2,-1) {1.1 Gobierno y contrato};
\node[anchor=east,align=right,text width=4.3cm] at (-.2,-2) {1.2 Plataforma y activos};
\node[anchor=east,align=right,text width=4.3cm] at (-.2,-3) {1.3 Seguridad y auditoría};
\node[anchor=east,align=right,text width=4.3cm] at (-.2,-4) {1.4 Núcleo operacional};
\node[anchor=east,align=right,text width=4.3cm] at (-.2,-5) {1.5 Núcleo de servicios};
\node[anchor=east,align=right,text width=4.3cm] at (-.2,-6) {1.6 Núcleo administrativo};
\node[anchor=east,align=right,text width=4.3cm] at (-.2,-7) {1.7 Canales e integración};
\node[anchor=east,align=right,text width=4.3cm] at (-.2,-8) {1.8 Migración e historia};
\node[anchor=east,align=right,text width=4.3cm] at (-.2,-9) {1.9 Innovaciones y pilotos};
\fill[SynBlue!55] (1,-9.22) rectangle (2,-8.78);
\fill[SynBlue!55] (2,-9.22) rectangle (3,-8.78);
\fill[SynBlue!55] (3,-9.22) rectangle (4,-8.78);
\fill[SynBlue!55] (4,-9.22) rectangle (5,-8.78);
\fill[SynBlue!55] (5,-9.22) rectangle (6,-8.78);
\fill[SynBlue!55] (6,-9.22) rectangle (7,-8.78);
\fill[SynBlue!55] (7,-9.22) rectangle (8,-8.78);
\fill[SynBlue!55] (8,-9.22) rectangle (9,-8.78);
\node[anchor=east,align=right,text width=4.3cm] at (-.2,-10) {1.10 Pruebas y correcciones};
\node[anchor=east,align=right,text width=4.3cm] at (-.2,-11) {1.11 Formación e implantación};
\node[anchor=east,align=right,text width=4.3cm] at (-.2,-12) {1.12 Servicio y salida};
\fill[SynBlue!55] (0,-12.22) rectangle (1,-11.78);
\fill[SynBlue!55] (1,-12.22) rectangle (2,-11.78);
\fill[SynBlue!55] (2,-12.22) rectangle (3,-11.78);
\fill[SynBlue!55] (3,-12.22) rectangle (4,-11.78);
\fill[SynBlue!55] (4,-12.22) rectangle (5,-11.78);
\fill[SynBlue!55] (5,-12.22) rectangle (6,-11.78);
\fill[SynBlue!55] (6,-12.22) rectangle (7,-11.78);
\fill[SynBlue!55] (7,-12.22) rectangle (8,-11.78);
\fill[SynBlue!55] (8,-12.22) rectangle (9,-11.78);
\fill[SynBlue!55] (9,-12.22) rectangle (10,-11.78);
\fill[SynBlue!55] (10,-12.22) rectangle (11,-11.78);
\fill[SynBlue!55] (11,-12.22) rectangle (12,-11.78);
\fill[SynBlue!55] (12,-12.22) rectangle (13,-11.78);
\fill[SynBlue!55] (13,-12.22) rectangle (14,-11.78);
\fill[SynBlue!55] (14,-12.22) rectangle (15,-11.78);
\fill[SynBlue!55] (15,-12.22) rectangle (16,-11.78);
\fill[SynBlue!55] (16,-12.22) rectangle (17,-11.78);
\fill[SynBlue!55] (17,-12.22) rectangle (18,-11.78);
\fill[SynBlue!55] (18,-12.22) rectangle (19,-11.78);
\fill[SynBlue!55] (19,-12.22) rectangle (20,-11.78);
\fill[SynBlue!55] (20,-12.22) rectangle (21,-11.78);
\fill[SynBlue!55] (21,-12.22) rectangle (22,-11.78);
\fill[SynBlue!55] (22,-12.22) rectangle (23,-11.78);
\fill[SynBlue!55] (23,-12.22) rectangle (24,-11.78);
\fill[SynBlue!55] (24,-12.22) rectangle (25,-11.78);
\fill[SynBlue!55] (25,-12.22) rectangle (26,-11.78);
\fill[SynBlue!55] (26,-12.22) rectangle (27,-11.78);
\fill[SynBlue!55] (27,-12.22) rectangle (28,-11.78);
\node[anchor=east,text width=4.3cm,align=right] at (-.2,-13) {Marcha blanca E1};
\node[anchor=east,text width=4.3cm,align=right] at (-.2,-14) {Marcha blanca E2};
\node[anchor=east,text width=4.3cm,align=right] at (-.2,-15) {Operación ambos alcances};
\fill[SynNavyLight!65] (0,-15.22) rectangle (28,-14.78);
\node[anchor=east,text width=4.3cm,align=right] at (-.2,-16) {Cobertura inicial E1/E2};
\end{tikzpicture}\caption{Carta Gantt agregada de los meses 29--56}\label{fig:sd7-gantt-2}
\FuenteFigura{Intervalos de red y asignaciones con esfuerzo UCP seleccionado; acta de marcha blanca E2 y corte requeridos al inicio de 21, bajo el supuesto adoptado de recepción continua y firma expresa programada.}
\end{figure}
Las filas permiten contrastar la continuidad de las cuentas con las ventanas contractuales. Los hitos identifican recepciones; la extensión de una barra no modifica las precedencias ni habilita trabajo durante una restricción operacional.
\end{landscape}\clearpage

